% RESEARCH OBJECTIVES

\chapter{Design and Planning}

In Chapter 3, students translate their research into a concrete engineering plan. This chapter directly addresses the syllabus requirements for "proposing multiple solutions and selecting the best possible solution" and "demonstrating project management principles".

\section{Proposed Solution}	
This section describes the "best possible solution" selected for the underline problem.

Core Approach: Detail the specific Machine Learning algorithm chosen (e.g., Random Forest, CNN, or LSTM) and why it is the most suitable for the identified societal problem.

Web Integration Strategy: Explain how the ML model will be served via the web (e.g., using Flask, Django, or FastAPI).

Innovation: Highlight the "innovativeness in solutions" as per the assessment criteria.

\section{System Architecture}	  
This provides a visual and technical roadmap of the "working model". Our system follows the three tier architecture. First tier consist of GUI, followed by Processing block and the Database.

\subsection{GUI (Graphical User Interface)}
In our project GUI deals with the interface for the user where the user will have to register and then login give the skill test and submit his resume in any format (Pdf, doc, docs, etc.)

\subsection{Processing block}
This is the block where the actual processing of our project is done. This block will connect the database and the GUI .i.e. it acts as a connector as well as communicator which connects the database and helps in transfer of data between the GUI and the database. Its main function will be to take input from the uploaded CV of the candidate.

\subsection{Databas:}
Database tier is used for the data storage. This contains all the data that is need for the processing and string data of the project. The data in this tier is related to the student information gathered from his/her CV.

\section{Design}

Include the UML Diagram as per requirement of your project
\subsection{Use case Diagram}

\subsection{Activity Diagram/ Sequence Diagram}

\subsection{Data Flow Diagram}

\section{Algorithm}
Explain the algorithms used in your project.
\subsection{K- Means Algorithm}
The Lloyd's algorithm, mostly known as k-means algorithm, is used to solve the k-means clustering problem and works as follows. First, decide the number of cluster:
Step 1: Initialize the center of the clusters
Step 2:.Attribute the closest cluster to each data point
Step 3: Set the position of each cluster to the mean of all data points belonging to that cluster Step 4 : Repeat steps 2-3 until convergence


\section{Component Selection}
This section documents the "system selection and cost analysis" required by the project.

Software Selection: Justify the choice of programming languages (e.g., Python), libraries (e.g., Scikit-learn), and web frameworks based on "engineering fundamentals".

Hardware Requirements: List the computational resources needed for training the model and hosting the web application.

Cost Analysis: Provide a detailed breakdown of costs (if any) to ensure "cost effectiveness," a key assessment criterion.

\section{Implementation Plan}
Weekly activity breakdown using a Gantt, PERT, or CPM chart.
Students must demonstrate "project management principles" by organizing their work into a schedule.

Weekly Activity Breakdown: Each group must provide an "implementation plan in the form of Gantt/PERT/CPM chart".

Milestones: Map key activities—such as "identification of need," "building prototype," and "testing"—to specific weeks of the semester.

Self-Learning Progress: Align the plan with the "log book" entries to record weekly work progress for verification by the project guide.

